% !TEX encoding = UTF-8 Unicode
\section{Realisierung}
\subsection{SDK Probleme}
%string längen
%Parameter::setValue( VstNumber v ) avoid NaN. problems with serialize and deserialize issue #113
%midi receive
%getLocation
%fruity
%Cant add vst-directories in VSTForx configuration dialog
Da VSTForx sowohl ein Plugin, als auch ein Host für Plugins ist, existieren gleich zwei potentielle Fehlerquellen. Zum einen muss sichergestellt werden, dass jede Host-Anwendung VSTForx Laden und Ausführen kann. Zum anderen sollte auch VSTForx in der Lage sein, jedes Plugin zu Laden und Auszuführen. Unglücklicherweise ist eine große schwäche des VST2.x-SKDs, die schlechte Dokumentation. So sind viele Funktionen ungenügend oder gar nicht beschrieben. Auch bleibt ungeklärt, in welcher Reihenfolge Operationen ausgeführt werden(siehe Abschnitt: Initalisierungsreihenfolge). Da es eine große Menge an Plugins und Host-Anwendungen gibt, führte dies zu einer Reihe an Problemen, die erst durch viel ausprobieren ermittelt und behoben werden konnten. 
\paragraph{Initalisierungsreihenfolge}
Wie bereits in Kapitel \ref{kapAnforderungen} beschrieben, existiert zur Kommunikation zwischen Host und Plugin lediglich eine Methode: \funktion{dispatcher()}.

// reihenfolge wichtig! ( ueber debugger ermittelt )
// setze samplerate  
plug.aEff->dispatcher ( plug.aEff, effSetSampleRate, 0, 0, 0, plug.hostInfo->getSampleRate() );
// setze blockSize  
plug.aEff->dispatcher ( plug.aEff, effSetBlockSize, 0, plug.hostInfo->getBlockSize(), 0, 0 );
// open
plug.aEff->dispatcher ( plug.aEff, effOpen, 0, 0, 0, 0.0 );
aEff->dispatcher ( aEff, effMainsChanged, 0, 1, 0, 0 ); // turn on
aEff->dispatcher ( aEff, effMainsChanged, 0, 0, 0, 0 ); // turn off
aEff->dispatcher ( aEff, effMainsChanged, 0, 1, 0, 0 ); // turn on


\paragraph{String längen}
\paragraph{NaN Probleme}
\paragraph{Kein Midi-Empfang}
\paragraph{Wo bin ich?}
\paragraph{Verzeichnisauswahl}

%%%%%%%%%%%%%%%%%%%%%%%%%%%%%%%%%%%%%%%%%%%%%%%%%%%%%%%%%%%%%%%%%%%%
\subsection{Persistenz}
%Serialisierungs problem : ein prozess
%%%%%%%%%%%%%%%%%%%%%%%%%%%%%%%%%%%%%%%%%%%%%%%%%%%%%%%%%%%%%%%%%%%%
\subsection{Die Plugindatenbank}
\paragraph{Testen}
\paragraph{Probleme}
%%%%%%%%%%%%%%%%%%%%%%%%%%%%%%%%%%%%%%%%%%%%%%%%%%%%%%%%%%%%%%%%%%%%
\subsection{Die Benutzerschnitstelle}
\paragraph{Probleme}
